% Section 8 -- related work, half a page. Every citation kept with its
% differentiating claim; the paper-by-paper elaboration is gone.

\section{Related work}
\label{sec:related}

\citet{li2026network} give the multilayer analysis of mathlib4 this work
consumes (the graph, the explicit/synthesized distinction, the two-layer hub
structure) and pose structural drift from machine-generated proofs as future
work; we extend their human baseline to a machine population rather than building a
baseline.

Premise retrieval asks which declarations should close a given goal and scores
a ranking against the premises a human cited. \citet{leandojo} supply that
setting and its annotations, and \citet{cslibpremise} establish source-visible
extraction for it, the precedent for our instrument. We run the relation
backwards, asking what the selection of already-verified proofs looks like in
aggregate against a library, with no right answer per goal. A retriever can
score well on every goal it is given while the population it serves still
occupies a small region of the library.

Three lines measure past correctness and each differs in its unit.
\citet{theorembench} score partial progress \emph{within} a theorem;
\citet{proofrank} characterize a proof text through proxies and pairwise
comparison, needing a preference model where ours needs none;
\citet{verisoftbench} correlate success with transitive dependency closure
size, a magnitude where coverage is a location, in a domain where the graph is
a build artifact. Our corpora are the released \citet{goedelprover} proof set
and mathlib itself; \citet{minif2f} appears for what it lacks, the human Lean~4
proofs that would permit a same-problem comparison.
